\documentclass{ctexart}
\ctexset{
    section = {
        titleformat = \raggedright,
        name = {,、},
        number = \chinese{section}
    }
}
\usepackage[a4paper,top=2.54cm,bottom=2.54cm,left=3.18cm,right=3.18cm]{geometry}
\usepackage{amsmath,amsfonts,amssymb,amsthm}
\usepackage{tikz}
\usepackage{multirow}
\usepackage{array}
\usepackage{fancyhdr}
\usepackage{lastpage}

\pagestyle{fancy}
\fancyhf{}
\cfoot{【第 \thepage 页~~共 \pageref{LastPage} 页】}
\renewcommand{\headrulewidth}{0pt}
\renewcommand{\labelenumii}{\Alph{enumii}.}

\begin{document}

\begin{center}
    【第 \thepage 页~~共 \pageref{LastPage} 页】
\end{center}

\begin{center}
\textbf{（主干课程）}

\textbf{（非主干课程）}

\textbf{（主干课程/非主干课程）}
\end{center}
\vspace{1em}
\begin{center}
\textbf{\LARGE 江西财经大学}\par
\vspace{8pt}
\textbf{\LARGE 2021--2022学年第一学期期末考试试卷}
\end{center}

\begin{center}
\begin{tabular}{m{0.33\textwidth} m{0.25\textwidth} m{0.28\textwidth}}
     试卷代码：1004203463A &授课课时：48 &考试用时：110分钟 \\
     课程名称：\textbf{金融科技概论} 
     & \multicolumn{2}{l}{适用对象：}\\
     试卷命题人：\underline{Yanfeng Wu} &           &试卷审核人：\underline{XXX}
\end{tabular}
\end{center}
\hrule height 1pt
%
oindent
%\rule{\textwidth}{1pt}

\section{选择题（共20分，每题4分）}
\begin{enumerate}
    \item 四选一（ ）
    \begin{enumerate}
        \item 1
        \item 2
        \item 3
        \item 4
    \end{enumerate}
    
    \item 四选一（ ）
    \begin{enumerate}
        \item 1
        \item 2
        \item 3
        \item 4
    \end{enumerate}
\end{enumerate}

\section{概念或名词解释（共20分，每题4分）}
\begin{enumerate}
    \item 中国
    \item China
    \item Jiangxi
    \item JXUFE
    \item Nanchang
    \item 中国
    \item China
    \item Jiangxi
    \item JXUFE
    \item Nanchang
\end{enumerate}

\section{计算题（共30分，每题15分）}
请证明勾股定理：
$$
a^2 + b^2 = c^2.
$$
%
ewpage



\section{证明或推导题（共30分，每题15分）}

\end{document}
